\documentclass[10pt]{article}
\usepackage[T1]{fontenc}
\usepackage[utf8]{inputenc}
\usepackage[french]{babel}
\usepackage{amsmath,amssymb,mathrsfs,mathpazo}
\usepackage{geometry}
\geometry{a4paper,margin=2cm}
\usepackage{graphicx,array,tabularx,tabularray,colortbl,multirow,multicol,booktabs,enumitem,pifont,ulem,fancybox,diagbox,cancel}
\usepackage{pstricks,pstricks-add,pst-plot,pst-tree,pst-node,pst-text,pst-func,pst-math,pst-bspline,pst-3dplot,pst-eucl,pst-solides3d}
\let\pscalculate\undefined
\usepackage{pst-bezier}
\usepackage{tikz,pgfplots,tkz-tab}
\usetikzlibrary{arrows,arrows.meta,patterns,calc,positioning,trees,decorations.markings,shapes.geometric,shapes.misc,shapes.arrows,intersections,decorations.pathmorphing,decorations.pathreplacing,backgrounds,shadows,babel}
\pgfplotsset{compat=1.18}
\usepgfplotslibrary{fillbetween}
\usepackage[np]{numprint}
\usepackage{listings,makecell,pas-tableur,twoopt}
\let\barre\undefined
\usepackage[pstricks]{tablvar}
\definecolor{darkgreen}{RGB}{0,100,0}
\providecommand{\red}{\color{red}}
\providecommand{\blue}{\color{blue}}
\providecommand{\green}{\color{green}}
\providecommand{\black}{\color{black}}
\providecommand{\white}{\color{white}}
\providecommand{\gray}{\color{gray}}
\providecommand{\cyan}{\color{cyan}}
\providecommand{\magenta}{\color{magenta}}
\providecommand{\yellow}{\color{yellow}}
\providecommand{\violet}{\color{violet}}
\providecommand{\orange}{\color{orange}}
\providecommand{\euro}{\texteuro}
\providecommand{\EUR}{\texteuro}
\providecommand{\ieme}{\textsuperscript{e}}
\providecommand{\degres}{\ensuremath{^\circ}}
\providecommand{\degree}{\ensuremath{^\circ}}
\providecommand{\vv}[1]{\overrightarrow{#1}}
\providecommand{\Coord}[1]{#1}
\providecommand{\decosix}{}
\providecommand{\starredbullet}{\ensuremath{\bullet}}
\pagestyle{empty}
\setlength{\parindent}{0pt}
\usepackage{adjustbox}\DeclareUnicodeCharacter{2236}{:}\DeclareUnicodeCharacter{2190}{\ensuremath{\leftarrow}}
\providecommand{\R}{\mathbb{R}}
\providecommand{\N}{\mathbb{N}}
\providecommand{\Z}{\mathbb{Z}}
\providecommand{\Q}{\mathbb{Q}}
\providecommand{\D}{\mathbb{D}}
\providecommand{\C}{\mathbb{C}}
\providecommand{\e}{{\mathrm{e}}}
\providecommand{\vect}[1]{\overrightarrow{#1}}
\providecommand{\vectt}[1]{\overrightarrow{\mathrm{#1}}}
\providecommand{\vvt}[1]{\overrightarrow{\mathrm{#1}}}
\providecommand{\np}[2][]{{#2~#1}}
\providecommand{\vv}[1]{{\overrightarrow{#1}}}
\renewcommand{\d}{{\mathrm{d}}}
\providecommand{\degres}{{^\circ}}
\providecommand{\degre}{{^\circ}}
\providecommand{\blue}{\color{blue}}
\providecommand{\red}{\color{red}}
\providecommand{\enskip}{\quad}
\providecommand{\llbracket}{\lbrack\!\lbrack}
\providecommand{\rrbracket}{\rbrack\!\rbrack}
\providecommand{\euro}{{\text{€}}}
\providecommand{\no}{{\text{n°}}}
\providecommand{\interval}[2]{{[#1;#2]}}
\providecommand{\intervFO}[2]{{[#1;#2[}}
\providecommand{\dotfill}{\ldots}

\providecommand{\ds}{\displaystyle}
\providecommand{\textding}[1]{#1}
\providecommand{\Oij}{\(\mathrm{O};\vec{\imath},\vec{\jmath}\)}
\providecommand{\Oijk}{\(\mathrm{O};\vec{\imath},\vec{\jmath},\vec{k}\)}
\providecommand{\Ouv}{\(\mathrm{O};\vec{u},\vec{v}\)}
\providecommand{\intervFF}[2]{\ensuremath{[#1;#2]}}
\renewcommand{\euro}{\ifmmode\text{\texteuro}\else\texteuro\fi}
\DeclareUnicodeCharacter{00B0}{\ifmmode^{\circ}\else\textdegree\fi}
\DeclareUnicodeCharacter{20AC}{\ifmmode\text{\texteuro}\else\texteuro\fi}
% Énoncés : Éducation nationale. Transcription et corrigés : APMEP, Denis Vergès et contributeurs. Découpage et classement : CoopMaths, Sébastien Lozano, Éric Elter et contributeurs. Adaptation de la présentation : Maths Explorer.
% CC BY-SA : https://coopmaths.fr/www/about/
% https://coopmaths.fr/alea/static/bac/2022/tex/bac_2022_05_sujet1_etrangers_3_cor.tex
\begin{document}
\noindent\textit{APMEP / CoopMaths. CC BY-SA. Présentation adaptée par Maths Explorer.}\par\medskip
\bigskip

\textbf{Partie A :}

\medskip

Soit $h$ la fonction définie sur $\R$ par

\[h(x) = \text{e}^x - x\]

\smallskip

\begin{enumerate}
\item 

$\bullet~~$$\displaystyle\lim_{x \to - \infty} \text{e}^x = 0$ et $\displaystyle\lim_{x \to - \infty} - x = + \infty$, donc par somme de limites : $\displaystyle\lim_{x \to - \infty} h(x) = + \infty$.

$\bullet~~$Pour $x \ne 0$, \: $h(x) = x\left(\dfrac{\text{e}^x}{x} - 1\right)$.

Or $\displaystyle\lim_{x \to + \infty} \dfrac{\text{e}^x}{x} = + \infty$, d'où $\displaystyle\lim_{x \to + \infty}\dfrac{\text{e}^x}{x} - 1 = + \infty$ et par produit de limites :

$\displaystyle\lim_{x \to + \infty} h(x) = + \infty$ ; 

\item 

Somme de fonctions dérivables sur $\R$, $h$ est dérivable sur $\R$ et sur cet intervalle :

$h'(x) =  \text{e}^x - 1$.

$\bullet~~$$h'(x) > 0 \iff \text{e}^x - 1 > 0 \iff \text{e}^x > 1 \iff \text{e}^x > \text{e}^0 \iff x > 0$ ;

$\bullet~~$$h'(x) < 0 \iff \text{e}^x - 1 < 0 \iff \text{e}^x < 1 \iff \text{e}^x < \text{e}^0 \iff x < 0$.

La fonction $h$ est donc décroissante sur $\R_{-}$ de plus l'infini à $h(0) = \text{e}^0 - 0 = 1$, puis croissante sur $\R_{+}$ de $h(0) = 1$ à plus l'infini.
\item 


$h$ étant croissante sur $\R_{+}$, \: $0 < a < b \Rightarrow h(0) < h(a) < h(b)$.

Comme $h(0) = 1$, on a donc $1 < h(a) < h(b) \Rightarrow h(a) - h(b) < 0$.
\end{enumerate}

\bigskip

\textbf{Partie B :}

\medskip

Soit $f$ la fonction définie sur $\R$ par

\[f(x) = \text{e}^x\]

On note $\mathcal{C}_f$ sa courbe représentative dans un repère \Oij.

\medskip

\begin{enumerate}
\item 
On a $M(x~;~y) \in T \iff y - f(0) = f'(0)(x - 0)$.

On a $f'(x) = f(x) = \text{e}^x$ d'où $f'(0) = \text{e}^0  = 1$ et $f(0) = \text{e}^0  = 1$.

$M(x~;~y) \in T \iff y - 1 = 1 (x - 0) \iff y = x + 1$.
\end{enumerate}

Dans la suite de l'exercice on s'intéresse à l'écart entre $T$ et $\mathcal{C}_f$ au voisinage de $0$. 

Cet écart est défini comme la différence des ordonnées des points de $T$ et $\mathcal{C}_f$ de même abscisse.

On s'intéresse aux points d'abscisse $\dfrac{1}{n}$, avec $n$ entier naturel non nul.

On considère alors la suite $\left(u_n\right)$ définie pour tout entier naturel non nul $n$ par : 

\[u_n = \text{exp} \left(\dfrac{1}{n}\right) - \dfrac{1}{n} - 1\]

\begin{enumerate}[resume]
\item 
On a $\displaystyle\lim_{n \to + \infty} \dfrac{1}{n} = 0$, d'où $\displaystyle\lim_{n \to + \infty} \text{e}^{\frac1n} = \text{e}^0 = 1$ et enfin $\displaystyle\lim_{n \to + \infty} u_n = 1 - 0 - 1 = 0$.
\item 
	\begin{enumerate}
		\item 
		



On a pour tout $n \in \N$, \: $u_{n+1} - u_n = \text{exp} \left(\dfrac{1}{n+1}\right) - \dfrac{1}{n+1} - 1 - \left(\text{exp} \left(\dfrac{1}{n}\right) - \dfrac{1}{n} - 1 \right) =$

$ \text{exp} \left(\dfrac{1}{n+1}\right) - \dfrac{1}{n+1} - 1 - \text{exp} \left(\dfrac{1}{n}\right) + \dfrac{1}{n} + 1 = \text{exp} \left(\dfrac{1}{n+1}\right) - \dfrac{1}{n+1}  - \text{exp} \left(\dfrac{1}{n}\right) + \dfrac{1}{n}  =$

$\text{exp} \left(\dfrac{1}{n+1}\right) - \dfrac{1}{n+1}  - \left[\text{exp} \left(\dfrac{1}{n}\right) - \dfrac{1}{n}\right] =  h\left(\dfrac{1}{n + 1}\right) - h\left(\dfrac{1}{n}\right)$.
		\item 
		
On a $n$ non nul, donc $0 < n < n + 1 \Rightarrow 0< \dfrac{1}{n+1} < \dfrac{1}{n} \Rightarrow h\left(\dfrac{1}{n+1}\right) < h\left(\dfrac{1}{n}\right) \Rightarrow$
		
		$ h\left(\dfrac{1}{n + 1}\right) - h\left(\dfrac{1}{n}\right) < 0$.
		
		Finalement pour $n$ non nul, $u_{n+1} - u_n < 0$ ce qui montre que la suite $\left(u_n\right)$ est décroissante.
	\end{enumerate}
\item 


















Le tableau montre que $u_8 \approx 0,008 < 10^{-2}$.

Donc pour $x = \dfrac{1}{8}$, l'écart entre la courbe et la tangente est inférieur à un centième.

\end{enumerate}

\bigskip

\end{document}
